\begin{objectivebox}[이 장에서 붙잡을 질문]
논리값 \(0\)과 \(1\)을 숫자의 자리로 읽으면, 같은 진리표로 여러 자리의 덧셈과 뺄셈까지 만들 수 있을까?
\end{objectivebox}

\section{같은 0과 1에 자리가 생긴다}

앞 장까지 \(0\)과 \(1\)은 꺼짐과 켜짐, 거짓과 참을 나타내는 논리값이었다. 이제 같은 두 기호를 수를 적는 숫자로 읽어 보자. 영어 \emph{binary digit}, 곧 이진 숫자를 줄인 말이 \keyterm{비트}(bit)다. 비트 하나만 놓고 보면 \(0\) 또는 \(1\)이라는 두 경우뿐이다. 그러나 여러 비트를 줄지어 놓고 각 위치에 서로 다른 크기를 주면 하나의 수가 된다.

십진수 \(507\)에서 오른쪽부터 \(1,10,100\)의 자리가 놓이는 것처럼, 이진수에서는 오른쪽부터 \(1,2,4,8\)의 자리가 놓인다. 각 자리는 오른쪽으로 갈 때마다 2로 나뉘고 왼쪽으로 갈 때마다 2배가 된다. 이 위치별 크기를 \keyterm{자릿값}(place value)이라고 한다.

예를 들어 네 비트 \(0101_2\)는 다음과 같이 읽는다. 아래첨자 \(2\)는 이 수를 십진수가 아니라 이진수로 적었다는 표시다.

\[
  0101_2
  =0\times8+1\times4+0\times2+1\times1
  =5.
\]

비트가 논리값이라는 사실은 사라지지 않았다. 각 자리는 여전히 \(0\) 또는 \(1\)이다. 달라진 것은 위치마다 \(8,4,2,1\)이라는 무게를 붙였다는 점이다. 논리 회로는 비트열을 처리하고, 사람은 정해 둔 표현 규약에 따라 그 비트열을 수로 해석한다\parencite{mit6004information}.

네 비트로 만들 수 있는 서로 다른 비트열은 \(2^4=16\)개다. 모든 자리가 \(0\)인 \(0000_2\)를 0으로, 모든 자리가 \(1\)인 \(1111_2\)를

\[
  8+4+2+1=15
\]

로 읽으면 0부터 15까지 열여섯 수를 나타낼 수 있다. 음수를 쓰지 않는 이 해석을 \emph{무부호}(unsigned) 정수 표현이라고 한다. 일반적으로 \(n\)비트 무부호 정수의 범위는 \(0\)부터 \(2^n-1\)까지다.

여기서 비트 폭은 단순한 표기상의 장식이 아니다. 네 비트 값 하나를 저장하려면 네 비트를 보존할 공간이 필요하고, 한 번에 옮기려면 그만큼의 신호선이나 여러 번의 이동이 필요하다. 비트 폭을 늘리면 더 큰 수를 나타낼 수 있지만, 저장하고 옮겨야 할 양도 함께 늘어난다.

\begin{bookfigure}{같은 비트열을 서로 다른 수로 읽을 수 있다}{fig:ch03-two-faces-of-bits}
\begin{tikzpicture}[diagram]
  \path[use as bounding box] (0,0) rectangle (13.0,7.20);

  \node[diagram panel title] at (0.10,6.88) {회로가 다루는 비트열};
  \node[diagram note,anchor=east,text=black!62] at (4.20,6.88) {논리값};
  \foreach \x/\v in {4.70/1,5.70/0,6.70/0,7.70/0}{
    \node[concept box,minimum width=8mm,minimum height=9mm] at (\x,6.42) {\v};
  }
  \node[diagram note,anchor=west,text=black!62] at (8.25,6.42) {네 칸의 \(0\)과 \(1\)};

  \draw[concept dashed arrow] (6.20,5.88) -- (3.25,5.06);
  \draw[concept dashed arrow] (6.20,5.88) -- (9.75,5.06);
  \node[diagram note,align=center,text=black!62] at (6.45,5.27) {표현 규약이\\자릿값을 정한다};

  \draw[diagram panel] (0.05,1.20) rectangle (6.20,4.95);
  \node[diagram panel title,anchor=north] at (3.13,4.62) {4비트 무부호 정수};
  \node[diagram note,anchor=east,text=black!62] at (1.16,3.82) {자릿값};
  \foreach \x/\v in {1.72/8,2.72/4,3.72/2,4.72/1}{
    \node[concept emphasis,minimum width=8mm,minimum height=8mm] at (\x,3.82) {\v};
  }
  \node[diagram note,align=center] at (3.12,2.74)
    {\(1000_2=1\times8+0\times4+0\times2+0\times1=8\)};
  \node[concept box,minimum width=34mm,minimum height=8mm] at (3.12,1.78) {표현 범위 \(0\sim15\)};

  \draw[diagram panel,fill=black!4] (6.80,1.20) rectangle (12.95,4.95);
  \node[diagram panel title,anchor=north] at (9.87,4.62) {4비트 2의 보수};
  \node[diagram note,anchor=east,text=black!62] at (7.90,3.82) {자릿값};
  \foreach \x/\v in {8.46/{-8},9.46/4,10.46/2,11.46/1}{
    \node[concept emphasis,minimum width=8mm,minimum height=8mm] at (\x,3.82) {\v};
  }
  \node[diagram note,align=center] at (9.87,2.74)
    {\(1000_2=1\times(-8)+0\times4+0\times2+0\times1=-8\)};
  \node[concept box,minimum width=34mm,minimum height=8mm] at (9.87,1.78) {표현 범위 \(-8\sim7\)};

  \draw[gate group] (0.10,0.22) rectangle (12.92,0.86);
  \node[diagram footer] at (6.50,0.54)
    {비트열은 같고, 수의 해석은 다르다 · 회로 결과와 표현 규약을 구분한다};
\end{tikzpicture}
\end{bookfigure}


\figref{fig:ch03-two-faces-of-bits}의 위쪽 비트열은 하나지만, 아래쪽 수는 표현 규약에 따라 갈라진다. 오른쪽의 \(-8\) 해석은 잠시 뒤 2의 보수를 배운 다음 완성한다. 중요한 것은 비트열과 그 비트열이 뜻하는 수를 먼저 구분하는 일이다. 회로가 내놓은 \(1000\)이라는 네 비트와 그것을 8 또는 \(-8\)로 읽는 규칙은 같은 층의 대상이 아니다.

\section{한 자리 덧셈에는 두 출력이 필요하다}

이제 이진수 한 자리를 더해 보자. \(0+0=0\), \(0+1=1\), \(1+0=1\)까지는 결과가 한 자리 안에 들어간다. 그러나 \(1+1=2\)이고, 2를 이진수로 적으면 \(10_2\)다. 현재 자리에는 \(0\)을 쓰고 왼쪽의 다음 자리로 \(1\)을 넘겨야 한다.

따라서 한 자리 덧셈은 출력 하나로 충분하지 않다. 현재 자리에 남길 \emph{합} \(S\)와 다음 자리로 보낼 \emph{올림} \(C\)가 함께 필요하다. 두 입력 \(A,B\)를 받아 이 두 출력을 만드는 함수를 \keyterm{반가산기}(half adder)라고 한다. 영어 이름의 \emph{half}는 다음 자리에서 들어오는 올림까지는 아직 받지 않는다는 뜻이다.

\begin{center}
\small
\begin{tabular}{cccc}
\toprule
\sffamily\bfseries \(A\) & \sffamily\bfseries \(B\) & \sffamily\bfseries 합 \(S\) & \sffamily\bfseries 올림 \(C\) \\
\midrule
0 & 0 & 0 & 0 \\
0 & 1 & 1 & 0 \\
1 & 0 & 1 & 0 \\
1 & 1 & 0 & 1 \\
\bottomrule
\end{tabular}
\end{center}

합 \(S\)의 마지막 열은 2장에서 본 XOR의 \(0,1,1,0\)과 같다. 두 입력이 서로 다를 때만 현재 자리에 \(1\)이 남는다. 올림 \(C\)의 마지막 열은 1장에서 본 AND의 \(0,0,0,1\)과 같다. 두 입력이 모두 \(1\)일 때만 다음 자리로 \(1\)을 보낸다. 논리 함수 두 개를 나란히 놓자 한 자리 덧셈기가 되었다.

\section{들어온 올림까지 더하면 전가산기가 된다}

여러 자리 수에서는 한 자리만 따로 계산할 수 없다. 바로 오른쪽의 낮은 자리에서 올림이 넘어올 수 있기 때문이다. 두 입력 \(A,B\)에 \emph{들어온 올림} \(C_{\mathrm{in}}\)(carry-in)까지 받아 합 \(S\)와 \emph{나가는 올림} \(C_{\mathrm{out}}\)(carry-out)을 만드는 함수를 \keyterm{전가산기}(full adder)라고 한다.

입력이 세 개이므로 가능한 조합은 \(2^3=8\)개다. 세 비트 가운데 \(1\)이 하나면 합은 \(01_2\), 두 개면 \(10_2\), 세 개면 \(11_2\)다. 이 규칙을 빠짐없이 적으면 다음 표가 된다.

\begin{center}
\small
\begin{tabular}{ccccc}
\toprule
\sffamily\bfseries \(A\) & \sffamily\bfseries \(B\) & \sffamily\bfseries \(C_{\mathrm{in}}\) & \sffamily\bfseries \(S\) & \sffamily\bfseries \(C_{\mathrm{out}}\) \\
\midrule
0 & 0 & 0 & 0 & 0 \\
0 & 0 & 1 & 1 & 0 \\
0 & 1 & 0 & 1 & 0 \\
0 & 1 & 1 & 0 & 1 \\
1 & 0 & 0 & 1 & 0 \\
1 & 0 & 1 & 0 & 1 \\
1 & 1 & 0 & 0 & 1 \\
1 & 1 & 1 & 1 & 1 \\
\bottomrule
\end{tabular}
\end{center}

표 전체를 행마다 외울 필요는 없다. \(S\)는 세 입력의 합에서 \(1\)의 자리에 남는 비트이고, \(C_{\mathrm{out}}\)은 \(2\)의 자리에 놓이는 비트다. 예를 들어 \(A=0\), \(B=1\), \(C_{\mathrm{in}}=1\)이면 합계가 2이므로 \(10_2\), 곧 \(S=0\), \(C_{\mathrm{out}}=1\)이다. 세 입력이 모두 \(1\)이면 합계가 3이므로 \(11_2\), 곧 두 출력이 모두 \(1\)이다.

전가산기는 완전히 새로운 마법 상자가 아니다. 먼저 \(A\)와 \(B\)를 반가산기로 더하고, 그 합에 들어온 올림을 두 번째 반가산기로 더한 뒤, 두 곳에서 생길 수 있는 올림을 합칠 수 있다. 이미 만든 작은 불 함수를 계층적으로 조합해 반가산기, 전가산기, 여러 비트 덧셈기로 올라가는 교육용 구조가 바로 이 생각을 따른다\parencite{nand2tetrisproject2}.

이 장에서는 전가산기 안의 NAND 게이트 수나 가장 긴 경로를 세지 않는다. 지금 필요한 것은 입력 세 비트에 대해 두 출력이 정확히 무엇이어야 하는지 정하는 명세다. 같은 전가산기 함수를 어떤 회로로 구현할지는 다시 여러 선택을 허용한다.

\section{한 자리 표 네 개를 옆으로 연결한다}

이제 \(0101_2\), 곧 5와 \(0011_2\), 곧 3을 더하자. 오른쪽 끝을 0번 자리라고 하고 왼쪽으로 \(1,2,3\)번을 붙인다. \(A_i\)와 \(B_i\)는 \(i\)번 자리의 입력, \(S_i\)는 그 자리의 합이다. 가장 낮은 자리로 처음 들어오는 올림은 없으므로 \(c_0=0\)으로 둔다. 자리 \(i\)에서 나간 올림 \(c_{i+1}\)은 바로 왼쪽 자리로 들어간다.

\[
\begin{array}{r}
  0101_2\\[-1mm]
{}+0011_2\\
\hline
  1000_2
\end{array}
\]

낮은 자리부터 네 번의 전가산기 계산을 끝까지 적으면 다음과 같다.

\begin{center}
\small
\begin{tabular}{cccccc}
\toprule
\sffamily\bfseries 자리 \(i\) & \sffamily\bfseries \(A_i\) & \sffamily\bfseries \(B_i\) & \sffamily\bfseries 들어온 올림 \(c_i\) & \sffamily\bfseries 합 \(S_i\) & \sffamily\bfseries 나가는 올림 \(c_{i+1}\) \\
\midrule
0 & 1 & 1 & 0 & 0 & 1 \\
1 & 0 & 1 & 1 & 0 & 1 \\
2 & 1 & 0 & 1 & 0 & 1 \\
3 & 0 & 0 & 1 & 1 & 0 \\
\bottomrule
\end{tabular}
\end{center}

0번 자리에서는 \(1+1+0=10_2\)이므로 \(S_0=0\), \(c_1=1\)이다. 1번 자리에서는 \(0+1+1=10_2\)이므로 다시 \(S_1=0\), \(c_2=1\)이다. 2번 자리도 \(1+0+1=10_2\)이고, 3번 자리에서는 \(0+0+1=1\)이므로 \(S_3=1\), 마지막 올림 \(c_4=0\)이다. 네 합 비트를 높은 자리부터 모으면 \(S_3S_2S_1S_0=1000_2\)이고, 무부호로 읽으면 8이다.

\begin{bookfigure}{한 자리 덧셈 상자 네 개가 4비트 합을 만든다}{fig:ch03-four-bit-addition}
\begin{tikzpicture}[diagram]
  \path[use as bounding box] (0,0) rectangle (13.0,6.45);

  \node[diagram note,anchor=west,text=black!62] at (0.20,6.08) {높은 자리};
  \node[diagram note,anchor=east,text=black!62] at (12.80,6.08) {낮은 자리};

  % adder:full
  \node[concept box,minimum width=25mm,minimum height=27mm] (fa3) at (1.75,4.15)
    {\textbf{전가산기 3}\\[1mm]\(A_3=0,\ B_3=0\)\\\(c_3=1\)\\[1mm]\(S_3=1\)};
  % adder:full
  \node[concept box,minimum width=25mm,minimum height=27mm] (fa2) at (4.90,4.15)
    {\textbf{전가산기 2}\\[1mm]\(A_2=1,\ B_2=0\)\\\(c_2=1\)\\[1mm]\(S_2=0\)};
  % adder:full
  \node[concept box,minimum width=25mm,minimum height=27mm] (fa1) at (8.05,4.15)
    {\textbf{전가산기 1}\\[1mm]\(A_1=0,\ B_1=1\)\\\(c_1=1\)\\[1mm]\(S_1=0\)};
  % adder:full
  \node[concept emphasis,minimum width=25mm,minimum height=27mm] (fa0) at (11.20,4.15)
    {\textbf{전가산기 0}\\[1mm]\(A_0=1,\ B_0=1\)\\\(c_0=0\)\\[1mm]\(S_0=0\)};

  \draw[path wire] (fa0.west) -- node[above,diagram note] {\(c_1=1\)} (fa1.east);
  \draw[path wire] (fa1.west) -- node[above,diagram note] {\(c_2=1\)} (fa2.east);
  \draw[path wire] (fa2.west) -- node[above,diagram note] {\(c_3=1\)} (fa3.east);
  \draw[path wire] (fa3.west) -- ++(-0.58,0) node[above,diagram note] {\(c_4=0\)};

  \foreach \nodeName in {fa3,fa2,fa1,fa0}{
    \draw[wire] (\nodeName.south) -- ++(0,-0.56);
  }

  \node[concept emphasis,minimum width=92mm,minimum height=11mm] at (6.50,1.78)
    {\(S_3S_2S_1S_0=1000_2\) \quad 무부호 해석: \(5+3=8\)};
  \node[diagram note,align=center,text=black!62] at (6.50,0.74)
    {이 그림은 값의 연결만 보여 준다 · 기다림과 내부 경로는 4장에서 센다};
\end{tikzpicture}
\end{bookfigure}


\figref{fig:ch03-four-bit-addition}은 같은 계산을 네 전가산기의 연결로 보여 준다. 오른쪽의 0번 자리에서 생긴 올림이 왼쪽 상자의 입력이 되고, 그 출력이 다시 다음 상자로 이어진다. 이 그림은 값이 어떤 순서로 연결되는지만 보여 준다. 각 전가산기의 내부 경로와 가장 왼쪽 결과가 언제 확정되는지는 4장에서 따로 센다.

\section{음수를 정하면 뺄셈도 덧셈이 된다}

네 비트에는 \(2^4=16\)개의 비트열만 있다. 이 열여섯 개를 0부터 15까지 모두 양수 쪽에 쓸 수도 있고, 일부를 음수에 나누어 줄 수도 있다. 현대의 고정 폭 정수에서 널리 쓰는 규약이 \keyterm{2의 보수}(two's complement)다\parencite{mit6004information}.

4비트 2의 보수에서는 맨 왼쪽 자리의 자릿값을 \(+8\)이 아니라 \(-8\)로 읽는다. 나머지 세 자리는 그대로 \(4,2,1\)이다. 따라서

\[
  1101_2=-8+4+0+1=-3
\]

이고, 나타낼 수 있는 범위는 \(-8\)부터 \(7\)까지다. 일반적인 \(n\)비트 2의 보수 범위는 \(-2^{n-1}\)부터 \(2^{n-1}-1\)까지다.

양수의 비트열에서 음수 표현을 구하는 손쉬운 절차도 있다. 모든 비트를 뒤집고 1을 더한다. 3의 4비트 표현 \(0011\)에 이 절차를 적용하면

\[
  0011 \longrightarrow 1100 \longrightarrow 1101
\]

이므로 \(1101\)을 \(-3\)으로 읽는다. 이제 \(5-3\)을 \(5+(-3)\)으로 바꾸어 같은 네 자리 덧셈기에 넣는다.

\[
\begin{array}{r@{\ }l}
  &0101 \quad (5)\\[-1mm]
{}+&1101 \quad (-3)\\
\hline
1\,&0010
\end{array}
\]

다섯 번째로 튀어나온 왼쪽의 \(1\)은 \(c_4\), 곧 carry-out이다. 네 비트로 계산하기로 한 고정 폭 덧셈에서는 이 추가 비트를 결과 밖으로 보내고 아래 네 비트 \(0010\)을 남긴다. \(0010_2=2\)이므로 원하는 \(5-3=2\)를 얻었다. 같은 덧셈기가 양수의 덧셈과 음수를 포함한 덧셈을 모두 처리했다. 달라진 것은 회로가 아니라 비트열을 읽는 표현 규약이다.

\begin{cautionbox}[carry-out과 signed overflow는 다르다]
같은 \(1000\)을 무부호로 읽으면 8이고, 4비트 2의 보수로 읽으면 \(-8\)이다. 따라서 앞의 \(0101+0011=1000\)은 무부호 계산 \(5+3=8\)로는 정확하지만, 부호 있는 4비트 계산으로는 양수 8을 표현할 수 없다. 4비트 2의 보수 범위가 \(-8\sim7\)이기 때문이다. 이것을 \emph{부호 있는 넘침}(signed overflow)이라고 한다.

반면 \(0101+1101=1\,0010\)에는 carry-out이 \(1\)이지만, 2의 보수 해석 \(5+(-3)=2\)는 범위 안에서 정확하다. 두 예는 carry-out과 signed overflow가 같은 신호가 아님을 보여 준다.

\begin{center}
\scriptsize
\begin{tabular}{
  >{\raggedright\arraybackslash}p{0.28\textwidth}
  >{\centering\arraybackslash}p{0.13\textwidth}
  >{\raggedright\arraybackslash}p{0.29\textwidth}
  >{\centering\arraybackslash}p{0.14\textwidth}}
\toprule
\sffamily\bfseries 4비트 덧셈 & \sffamily\bfseries carry-out & \sffamily\bfseries 부호 있는 해석 & \sffamily\bfseries signed overflow \\
\midrule
\(0101+0011=1000\) & 0 & \(5+3\)의 8을 표현하지 못함 & 1 \\
\(0101+1101=1\,0010\) & 1 & \(5+(-3)=2\)를 정확히 표현함 & 0 \\
\bottomrule
\end{tabular}
\end{center}

회로가 잘못 더한 것이 아니다. 회로는 정해진 비트 폭에서 같은 이진 덧셈을 수행했다. 어떤 수를 뜻한다고 약속했는지, 정확한 답이 그 표현 범위 안에 들어오는지가 따로 문제다. 이 장에서는 두 신호의 차이까지만 확인하고 상세 검출 회로는 만들지 않는다.
\end{cautionbox}

\section{같은 회로와 같은 비트열도 해석은 다를 수 있다}

지금까지의 층을 분리해 보자. 첫째, 비트열은 \(0101\), \(1000\), \(1101\)처럼 \(0\)과 \(1\)이 놓인 모양이다. 둘째, 표현 규약은 그 모양을 무부호 수 또는 2의 보수 정수로 읽는 방법이다. 셋째, 덧셈 회로는 각 자리에서 합과 올림을 계산해 새로운 비트열을 만든다.

이 셋을 섞으면 \(1000\)이 8이면서 \(-8\)이라는 말이 모순처럼 보인다. 그러나 같은 글자 ‘10’도 십진수로 읽으면 10이고 이진수로 읽으면 2인 것과 다르지 않다. 비트열 자체에는 부호가 붙어 있지 않다. 비트 폭과 표현 규약이 그 뜻을 정한다.

비트 폭을 넓히면 표현 범위가 커지고 넘침이 줄어든다. 그 대신 값 하나마다 더 많은 저장공간이 필요하고, 더 많은 비트를 옮겨야 하며, 덧셈기에 연결할 자리도 늘어난다. 범위는 정답을 담을 수 있는 조건이고, 저장 비트·배선·회로는 그 조건을 실현하는 물리적 비용이다. 두 종류의 숫자를 한 단위처럼 합치지 않는 이유다.

이 장에서는 네 전가산기를 연결해 \(5+3=8\)과 \(5-3=2\)를 완성했다. 하지만 연결에는 아직 풀지 않은 문제가 있다. 3번 자리는 2번 자리의 올림을 받고, 2번 자리는 1번 자리를, 1번 자리는 0번 자리를 기다린다. 비트 폭이 커질수록 이 기다림은 어떻게 늘어날까? 다음 장에서는 같은 덧셈의 답을 유지하면서 올림을 한 자리씩 전달하는 방법과 더 많은 논리를 써서 미리 계산하는 방법을 비교한다.

\section*{이번 장의 선택과 비용}
\addcontentsline{toc}{section}{이번 장의 선택과 비용}

\begin{center}
\small
\begin{tabular}{>{\raggedright\arraybackslash\sffamily\bfseries}p{0.29\textwidth} >{\raggedright\arraybackslash}p{0.58\textwidth}}
\toprule
보존한 기능·정확한 결과 & 무부호 \(0101+0011=1000\), 곧 \(5+3=8\); 2의 보수 \(0101+1101=1\,0010\)에서 네 비트 결과 \(5-3=2\) \\
\midrule
계산이 요구하는 양 & 전가산기 명세 8행, 네 자리의 합·올림 계산, 자리 사이를 잇는 올림 세 개와 마지막 carry-out \\
\midrule
기계가 제공할 자원 & 값 하나마다 4비트, 자리별 합·올림 기능 4개와 비트·올림을 잇는 연결 통로 \\
\midrule
유한 정밀도에서 달라질 수 있는 것 & 4비트 무부호 범위 \(0\sim15\), 2의 보수 범위 \(-8\sim7\); 같은 \(1000\)의 수 해석과 넘침 여부가 달라짐 \\
\midrule
설계 대가 & 비트 폭을 늘리면 표현 범위가 커지지만 저장 비트, 이동할 비트, 배선과 덧셈 자리도 늘어남 \\
\midrule
아직 해결되지 않은 제약 & 모든 자리가 바로 아래 자리의 올림을 기다릴 때 가장 왼쪽 결과까지 이어지는 의존 경로 \\
\bottomrule
\end{tabular}
\end{center}
